% Integrated FF.40; review copy only, do not input separately.
\begin{claim}{FF.40}{Proved}\label{clm:FFmetabeliantest}
Let \(I=G_{K_p}\), let \(I'=\overline{[I,I]}\),
let \(I''=\overline{[I',I']}\), and put
\[
                  B_p=\overline{\mathbb Q}^{\,I''}.
                                                        \tag{FF.40.1}
\]
This maximal metabelian extension of \(K_p\)
is a proper countable algebraic subfield of
\(\overline{\mathbb Q}\), with henselian
value group \(\mathbb Q\), residue field
\(\overline{\mathbb F}_p\), and
\(\mu_\infty\subset B_p\).
For every finite extension \(E/B_p\),
\[
             C_{\operatorname{Aut}(\overline{\mathbb Q})}(G_E)=1.
                                                        \tag{FF.40.2}
\]
It belongs to none of the sufficient field
families established in this section through
Claim~\ref{clm:FFnilpotentreconstruct}, including
the subfields of finite extensions of a
fixed bounded-nilpotency stage at any prime,
and to none of the negative centralizer-ghost
families. Moreover,
\[
 K_p^{{\rm ab},\times}\subset
                  \bigcap_{n\ge1}(B_p^\times)^{p^n}.
                                                        \tag{FF.40.3}
\]
Thus the exact \(\mu_\infty\) Kummer-kernel
input used in Claims~\ref{clm:FFcyclotomicreconstruct}--
\ref{clm:FFnilpotentreconstruct} fails here.
No reconstruction failure is asserted.
\end{claim}
\begin{proof}
The field \(B_p\) contains
\(K_p^{\rm ab}\), hence \(U_p\).
The valuation, residue and roots-of-unity
assertions follow. It is countable and
of transcendence degree zero because
it is algebraic over \(\mathbb Q\).

We verify the needed nonabelian groups
explicitly. For any finite number model
\(M/\mathbb Q_p\), wild inertia
\(P_M\) is free pro-\(p\) of
countably infinite rank. Choose eight
free generators and map them to
\(1+E_{i,i+1}\), \(1\le i\le8\),
in \(\operatorname{UT}_9(\mathbb F_p)\).
The words
\[
 u=[[x_1,x_2],[x_3,x_4]],\qquad
 v=[[x_5,x_6],[x_7,x_8]]
\]
belong to \(P_M''\), and map to
\(1+E_{1,5}\) and \(1+E_{5,9}\).
Their commutator maps to
\(1+E_{1,9}\ne1\). Thus
\(P_M''\) is nonabelian for every
prime, including \(p=2\).
In particular \(I''\ne1\), proving
that \(B_p\ne\overline{\mathbb Q}\).

For a finite extension \(E/B_p\),
choose a number field \(D\subset E\)
with \(E=DB_p\), and write
\(M=D_v\). Its absolute group is
\[
                 H=G_E=I''\cap G_M
                             \triangleleft G_M.
                                                        \tag{FF.40.4}
\]
It lies in \(P_M\) and contains
\(P_M''\), so is nonabelian.
Its centralizer in \(P_M\) is
trivial. Indeed, if \(a\ne1\)
centralized \(H\), the finitely
generated closed free pro-\(p\)
subgroup generated by \(a,u,v\)
would put two noncommuting elements
in the procyclic centralizer of
\(a\), contrary to the free-group
centralizer theorem used in
Claim~\ref{clm:FFcyclotomictest}.

The subgroup \(C_{G_M}(H)\) is
normal in \(G_M\), by (FF.40.4).
Its commutator with \(P_M\) lies
in \(C_{G_M}(H)\cap P_M=1\),
so it centralizes \(P_M\).
The fixed field of \(P_M\) is
a finite extension of \(T_p\).
It reconstructs by
Claim~\ref{clm:FFtamereconstruct};
the necessary condition in
Claim~\ref{clm:FF6} therefore makes
the centralizer of \(P_M\) in
\(\operatorname{Aut}(\overline{\mathbb Q})\)
trivial. Thus \(C_{G_M}(H)=1\).

Any global automorphism centralizing
\(H\) preserves its fixed field
\(E\). The rank-one henselian
valuation of \(E\) is unique:
two distinct such valuations would
make it algebraically closed, by
the two-valuation argument proved
in Claim~\ref{clm:FFstricttest},
whereas \(H\) is nonabelian.
Its unique prolongation to
\(\overline{\mathbb Q}\) is
therefore preserved. The global
centralizing automorphism belongs
to \(G_{\mathbb Q_p}\); a positive
power belongs to the open subgroup
\(G_M\) and is then trivial.
Torsion-freeness of \(G_{\mathbb Q_p}\),
verified in Claim~\ref{clm:FFcyclotomictest},
makes the automorphism itself
trivial. This proves (FF.40.2).

We next exclude the new nilpotent
families, including abstract field
embeddings. Let \(s=1\) for
odd \(p\), and \(s=2\) for
\(p=2\). For arbitrarily large
\(n\), the finite \(p\)-group
\[
 Q_n=\langle a,b:
       a^{p^n}=b^{p^n}=1,\quad
       bab^{-1}=a^{1+p^s}\rangle
                                                        \tag{FF.40.5}
\]
is a semidirect product of two
cyclic groups of order \(p^n\).
The displayed action has order
dividing \(p^n\), so the
presentation is consistent.
It is metabelian, and for every
fixed \(c\),
\[
           \gamma_{c+1}(Q_n)
                  =\langle a^{p^{sc}}\rangle,
            \qquad
           |\gamma_{c+1}(Q_n)|=p^{n-sc}
                   \quad(n>sc).
                                                        \tag{FF.40.6}
\]
Every finite \(p\)-group is a
quotient of \(I\), by the
Frattini lifting proof in
Claim~\ref{clm:FFnilpotentreconstruct}.
Since \(Q_n''=1\), these quotients
factor through \(I/I''\).
Consequently the image of
\(\gamma_{c+1}(I)\) in
\(I/I''\) is infinite, by
(FF.40.6).
If \(B_p\) were contained in
a finite extension \(L/K_{p,c}\),
the subgroup \(G_L\) would be
open in \(\gamma_{c+1}(I)\)
and contained in \(I''\).
That would make this image
finite, a contradiction.

Suppose instead that an abstract
embedding of \(B_p\) had image
in a finite extension \(L/K_{q,c}\).
Extend it to an automorphism
\(g\) of \(\overline{\mathbb Q}\).
The transported henselian valuation
of \(gB_p\) extends to a henselian
valuation on the algebraic extension
\(L\). Its residue characteristic
is \(p\); the given valuation
of \(L\) has residue characteristic
\(q\). Since \(L\) is not
algebraically closed, uniqueness
of a rank-one henselian valuation
forces equality of the valuations
and \(p=q\). Their unique
prolongations to the algebraic
closure then give \(g\in G_{\mathbb Q_p}\).
Conjugation by this group preserves
both \(I''\) and the closed
lower central series of \(I\).
Thus \(gB_p=B_p\), reducing to
the already excluded inclusion.

For completeness, the preceding
positive families are excluded
by their already verified criteria.
The field contains \(\mu_\infty\),
which is a nontrivial divisible
multiplicative subgroup, so it
cannot be a Kummer-faithful
sub-\(q\)-adic field for any
\(q\). The values of
\(\zeta_{q^m}-1\) have unbounded
\(q\)-denominators at every
residue-characteristic-\(q\)
place, excluding the bounded-wild
families at all primes.
Its transcendence degree zero
excludes exactly the families with an explicit
positive-transcendence hypothesis in
Claim~\ref{clm:FFclasses}(3)--(4) and the
finitely generated transcendental fields
of Claim~\ref{clm:FFfunction}.
We do not assert arbitrary subfield
closure for a proper normal envelope
over a transcendental finitely generated field.

The proper-normal-envelope families
are excluded by the same complete
two-valuation argument as in
Claim~\ref{clm:FFstricttest}.
Explicitly, for any number field
\(D\) in a finite extension
\(E/B_p\), choose
\(g\in G_D\) moving the selected
valuation. Such a \(g\) exists
because the decomposition group
of a finite place of a number
field is proper; the local-square,
global-nonsquare CRT construction
in that claim proves this directly.
The fields \(E\) and \(gE\)
then give two distinct rank-one
henselian valuations on their
compositum. They remain distinct
there because otherwise their
unique prolongations to the
algebraic closure would coincide.
Hence \(EgE=\overline{\mathbb Q}\).
The normal closure of \(E/D\)
is therefore \(\overline{\mathbb Q}\).
Together with the precise normal-closure
equivalence in Claim~\ref{clm:FFclasses}(2),
this excludes the formation subfield
families as stated there. It also
excludes containment in a
proper normal algebraic extension
of any number field, after
adjoining that number field if
necessary. It also excludes being
finite over such a proper normal
field: two conjugate copies of
\(G_E\) would then be open in
its normal Galois group but have
trivial intersection, forcing
the infinite group \(G_E\) to
be finite. Finally (FF.40.2)
excludes all the negative
centralizer-ghost families.

To prove (FF.40.3), take
\(a\in K_p^{{\rm ab},\times}\).
Over \(K_p^{\rm ab}\), adjoin
all compatible \(p\)-power
roots of every conjugate of
\(a\) under \(I\). There are
only finitely many such conjugates,
since \(a\) is algebraic over
\(K_p\). All roots of unity
are already in the base, so
this is an abelian pro-\(p\)
extension of \(K_p^{\rm ab}\),
stable under \(I\). Its Galois
group over \(K_p\) is an
extension of an abelian group
by an abelian group and hence
metabelian. It lies in \(B_p\).
In particular \(a\) has
compatible \(p\)-power roots
in \(B_p\), proving (FF.40.3).
The element \(a=p\) already
shows that the Kummer kernel
is larger than \(\mu_\infty\).
\end{proof}

The field \(B_p\), for example \(B_2\),
is consequently a specified remaining
test with the smallest possible
cardinality and transcendence degree
in characteristic zero. An
inclusion-minimal field is not
claimed. Whether \(\mathcal R(B_p)\)
holds remains \textbf{Open} in
this work. The failed exact
Kummer-kernel step identifies a
precise boundary of the preceding
proofs; it is not a topos map or
a refutation of centralizer
sufficiency.
